\documentclass[11pt,a4paper]{article}
\usepackage[utf8]{inputenc}
\usepackage{amsmath,amsfonts,amssymb,bm}
\usepackage{geometry}
\geometry{margin=2.5cm}

\title{Waveprobe Quadratic Unconstrained Binary Optimization (QUBO) Specification}
\author{Sovereign Stack Research}
\date{2026-04-17}

\begin{document}

\maketitle

\section{Introduction}

The \textbf{Waveprobe} is a local selection kernel designed to identify high-coherence state injection candidates within a search window $W$. It formalizes the transition from a diffuse search manifold to a discrete execution trace by evaluating the overlap energy between the active engram state and historical witness traces.

\section{The Waveprobe State}

Let $|\psi(x)\rangle \in \mathbb{C}^n$ be the complex-valued state vector representing the informatic properties of a byte sequence $x$. The state is composed of discrete basis states $|e_i\rangle$:
\begin{equation}
|\psi(x)\rangle = \sum_{i=1}^{k} c_i |e_i\rangle
\end{equation}
where $c_i \in \mathbb{C}$ are coefficients encoding local features. The state is typically normalized such that $\langle \psi | \psi \rangle = 1$.

\section{Projector and Local QUBO Formalism}

The Waveprobe selection kernel acts as a projection operator $\hat{P}$ constructed from the current active state $|\psi_{\text{curr}}\rangle$:
\begin{equation}
\hat{P} = |\psi_{\text{curr}}\rangle \langle \psi_{\text{curr}}|
\end{equation}

The \textbf{overlap energy} $E$ between the current state and a past state $|\psi_{\text{past}}\rangle$ is given by the quadratic form characteristic of a QUBO solver:
\begin{equation}
E(s) = \langle \psi_{\text{past}} | \hat{P} | \psi_{\text{past}} \rangle = |\langle \psi_{\text{curr}} | \psi_{\text{past}} \rangle|^2
\end{equation}

In matrix form, the local selection problem is represented as a QUBO matrix $Q$:
\begin{equation}
Q_{ij} = \bar{c}_i c_j
\end{equation}
where the objective is to maximize $x^\dagger Q x$ over the candidate states in the search window.

\section{Phase-Lock Coherence and Feature Fusion}

The magnitude of the overlap is governed by the \textbf{phase-lock coherence} $\phi(s, x)$. This scalar signal fuses three primary informatic features:
\begin{equation}
\phi(s, x) = w_e \phi_e + w_r \phi_r + w_d \phi_d
\end{equation}
with the canonical weights:
\begin{equation}
w_e = 0.4, \quad w_r = 0.3, \quad w_d = 0.3
\end{equation}
where:
\begin{itemize}
    \item $\phi_e$ is the normalized Shannon entropy component.
    \item $\phi_r$ is the byte repetition rate component.
    \item $\phi_d$ is the dictionary potential (4-gram density) component.
\end{itemize}

\section{Indefinite Causal Order and Superposition}

Consistent with the wave-particle duality of the engram, past states exist in a state of \textbf{indefinite causal order} before injection. The selection event $\arg\max E(s)$ retroactively collapses the history:
\begin{itemize}
    \item \textbf{Wave state}: Diffuse contribution to the torsion gradient.
    \item \textbf{Particle state}: Discrete, traceable anchor for compression.
\end{itemize}

The causal order observable $\mathcal{O}_{AB}$ distinguishes between classical matching and quantum-inspired advantage via the Bell-like bound:
\begin{equation}
|\langle \mathcal{O}_{AB} \rangle| \leq 2
\end{equation}
Violations of this bound ($>2$) confirm the non-classical nature of the Waveprobe selection.

\section{Regret-Blink Coupling}

The search process is synchronized with the \textbf{Gated Epoch Synchronizer (GES)}. The timing $\Delta t_{\text{blink}}$ is coupled to the regret magnitude $R_{\text{mag}}$:
\begin{equation}
\Delta t_{\text{blink}} = 500\text{ms} + 200\text{ms} \cdot R_{\text{mag}}
\end{equation}
where $R_{\text{mag}}$ measures the prediction surprise. The \textbf{decoherence time} $t_{\text{dec}} = 200\text{ms}$ defines the refractory window during which the torsion frame correction is computed before the next emission.

\section{Conservation and Totality}

The Waveprobe is \textbf{information-conservative}. The injection of a state $s_{\text{probe}}$ is permitted if and only if it does not increase the local bits-per-byte (BPB) cost:
\begin{equation}
\text{BPB}(x, s_{\text{probe}}) \leq \text{BPB}(x, s_{\text{local}})
\end{equation}
This ensures that the "spoiler" player (boundary pressure) cannot terminate the trace, preserving the topological integrity of the manifold.

\end{document}
